\section{The cubic resolvent and exterior multiplicities}
\label{sec:exterior}

The obstruction in Section~\ref{sec:rigidity} is detected by an exterior
vector concentrated in one pair of opposite coordinates.  We shall count
such vectors through their three signed opposite products.  After packet
normalization these products form a triple in a cubic field, and
concentration near one exterior plane confines
two of them to a small rectangle.  Our goal is a uniform count in that
rectangle.  We first define the count and reduce its multiplicities to local
densities.  We then construct the integral quadratic structure that will
control those densities in the next two sections.

\subsection{The three products and the counting problem}

Let \(\widetilde K\subset\R\) be the normal closure of \(K\), and let
\(G=\Gal(\widetilde K/\Q)\) act on the four ordered embeddings of \(K\).
This permutation action is primitive: a nontrivial block would give a field
strictly between \(\Q\) and \(K\).  A transitive subgroup of \(S_4\) has
order \(4,8,12\), or \(24\).  A subgroup of order \(4\) or \(8\) is contained
in a Sylow \(2\)-subgroup, which preserves a partition into two pairs.
It follows that
\[
G=\mathfrak A_4\quad\hbox{or}\quad G=S_4.
\]
Both groups act transitively on the six two-element subsets of
\(\{1,2,3,4\}\) and on the set of pair partitions
\[
\mathcal B=\{12|34,\ 13|24,\ 14|23\}.
\]
Define
\begin{equation}
 F=\left\{(z_b)_{b\in\mathcal B}\in\widetilde K^{\mathcal B}:
       g(z_b)=z_{gb}\ \text{for all }g\in G,\ b\in\mathcal B\right\}.
 \label{eq:ext-resolvent}
\end{equation}
Evaluation at any partition identifies this algebra with the fixed field
of its stabilizer.  Thus \(F\) is a totally real cubic field, and its three
real embeddings are the coordinates in Equation~\eqref{eq:ext-resolvent}.
In particular, \(\Tr_{F/\Q}z=\sum_b z_b\).  This is the classical cubic
resolvent from the action on pair partitions; for its relation to cubic
resolvent rings, see \cite[Section~2.3]{Bhargava2004III}.  We write
\(R_F=\OO_F\) and \(D_F=|\Disc(F)|\).

Put
\[
 V=\bigwedge\nolimits^2_{\Q}K,\qquad
 L_V=\bigwedge\nolimits^2_{\Z}\OO.
\]
For \(w=\sum_r x_r\wedge y_r\in V\), its embedding coordinates are
\[
 w_{ij}=\sum_r
 \bigl(\sigma_i(x_r)\sigma_j(y_r)-\sigma_j(x_r)\sigma_i(y_r)\bigr)
 \quad(i<j),\qquad w_{ji}=-w_{ij}.
\]
Choose a basis \(\beta_1,\ldots,\beta_4\) of \(\OO\) whose real embedding
matrix \(A_{\OO}=(\sigma_i(\beta_j))\) has positive determinant.  Write
\(\delta_{\OO}=\det A_{\OO}\in\widetilde K\).  Its selected real value is
\(X=q\sqrt D\), and
\begin{equation}
 \delta_{\OO}^{\,2}=\Disc(\OO)=X^2.
 \label{eq:ext-embedding-determinant}
\end{equation}
Define the quadratic expression
\begin{equation}
 \Psi(w)=\delta_{\OO}^{-1}
       (w_{12}w_{34},-w_{13}w_{24},w_{14}w_{23}).
 \label{eq:ext-product-map}
\end{equation}

\begin{lemma}[The product map]
\label{lem:ext-products}
The expression in Equation~\eqref{eq:ext-product-map} defines a rational
quadratic map \(\Psi:V\to F\).  In the exterior basis of \(L_V\) induced
by \((\beta_i)\), write \(w=\sum_{i<j}t_{ij}\beta_i\wedge\beta_j\).
Then
\begin{equation}
 Q_V(w):=\Tr_{F/\Q}\Psi(w)
       =t_{12}t_{34}-t_{13}t_{24}+t_{14}t_{23}.
 \label{eq:ext-integral-form}
\end{equation}
If \(w\in V\) is nonzero, all six \(w_{ij}\), and hence all three
components of \(\Psi(w)\), are nonzero.  For every extension \(k/\Q\),
\(a\in(K\otimes_\Q k)^\times\), \(c\in k\), and \(w\in V\otimes_\Q k\),
\begin{equation}
 \Psi\bigl(c(\bigwedge\nolimits^2 a)w\bigr)
       =c^2\Nm_{K\otimes k/k}(a)\Psi(w),
 \label{eq:ext-covariance}
\end{equation}
where \(\bigwedge^2a\) is the exterior square of multiplication by \(a\).
\end{lemma}

\begin{proof}
An element \(g\in G\) permutes the rows of \(A_{\OO}\), so it multiplies
\(\delta_{\OO}\) by the sign of that permutation.  Permuting the four
indices in the three signed products in
Equation~\eqref{eq:ext-product-map} permutes the partitions and introduces
the same common sign.  Division by \(\delta_{\OO}\) gives precisely the
equivariance in Equation~\eqref{eq:ext-resolvent}.  This is an identity of
quadratic polynomials in the rational coordinates of \(w\), and therefore
defines a rational quadratic map to \(F\).

The Pfaffian identity for the exterior square of a four-dimensional
linear map is
\[
 Q\bigl((\bigwedge\nolimits^2 A_{\OO})t\bigr)
       =\det(A_{\OO})Q(t),
 \qquad
 Q(t)=t_{12}t_{34}-t_{13}t_{24}+t_{14}t_{23}.
\]
It follows either by expanding the six minors or by evaluating
\(\frac12 t\wedge t\) in the top exterior power.  This proves
Equation~\eqref{eq:ext-integral-form}.  If a coordinate of a rational
\(w\) vanishes, Galois transitivity on the unordered pairs makes all six
coordinates vanish.  The exterior embedding map is injective, so \(w=0\).
Finally, multiplication by \(a\) multiplies \(w_{ij}\) by \(a_i a_j\)
after splitting, where \(a_i\) are its four embeddings.  Each opposite
product is multiplied by \(\prod_i a_i=\Nm(a)\), and the scalar \(c\)
contributes \(c^2\).  This proves Equation~\eqref{eq:ext-covariance}
after splitting and hence over \(k\).
\end{proof}

We relate these products to the packet model of Section~\ref{sec:packets}.
Let \(I\) be one of its maximal-order ideal-class representatives, let
\(L\subset I\) be one of its local labels, and let
\(v=(v_i)\in T_\infty\).  If \(w\in\bigwedge^2_{\Z}L\), the corresponding
exterior vector in the normalized lattice has coordinates
\[
 \widehat w_{ij}=(X\Nm I)^{-1/2}w_{ij}v_i v_j .
\]
Since \(\prod_i|v_i|=1\), the number
\(\epsilon(v)=\prod_i v_i\) is \(1\) or \(-1\).  Thus
\begin{equation}
 (\widehat w_{12}\widehat w_{34},
  -\widehat w_{13}\widehat w_{24},
  \widehat w_{14}\widehat w_{23})
   =\epsilon(v)\,z(I,w),\qquad
 z(I,w):=\frac{\Psi(w)}{\Nm I}.
 \label{eq:ext-signed-products}
\end{equation}
In particular \(Q(\widehat w)=\epsilon(v)\Tr z(I,w)\).  This records
exactly how the diagonal signs act on the products.

Lemma~\ref{lem:ext-products} and the normalization above show that every
lattice in a packet lies outside each \(\mathcal E_b\).  To prove that weak
limits give these events zero mass, we need quantitative control of exterior
vectors near \(V_b\); the count below supplies the arithmetic input for that
control.

For a rational prime \(p\), fix a generator \(I_p=\xi_pR_p\) as in the
packet model, and give \(R_p^\times\) its Haar probability.  For
\(0\ne w\in\bigwedge^2_{\Z}I\), define
\begin{equation}
 P_p(I,w)=\int_{R_p^\times}
 \mathbf1_{\bigwedge^2_{\Z_p}(\xi_p t\OO_p)}(w)\,d^\times t .
 \label{eq:ext-label-probability}
\end{equation}
Changing \(\xi_p\) by an \(R_p\)-unit leaves this probability unchanged.
For \(p\nmid q\), it is \(1\).  Hence \(\prod_p P_p(I,w)\) is a finite
product.  Multiplication by \(U^+\) preserves these probabilities by
translation of the local unit variables.  For \(e\in U^+\), the algebraic
norm \(\Nm_{K/\Q}(e)=1\), so Equation~\eqref{eq:ext-covariance} also
shows that \(z(I,w)\) is unchanged.  We use \(U^+\) to unfold the positive
real parameters on each choice of signs.
Its action on nonzero exterior vectors is free: if a totally positive
unit \(e\) fixes \(w\), then \(e_i e_j=1\) for every pair by
Lemma~\ref{lem:ext-products}; the four positive \(e_i\) must all be \(1\).

For \(z\in F\), define the nonnegative weighted multiplicity
\begin{equation}
 M(z)=
 \sum_{\substack{I,\;
    w\in(\bigwedge^2_{\Z}I\setminus\{0\})/U^+\\
    z(I,w)=z}}
       \prod_p P_p(I,w).
 \label{eq:ext-multiplicity}
\end{equation}
The sum is initially allowed to have value \(+\infty\); the local bounds
below prove its finiteness on the nonzero trace slices used here.
It depends on \(K,\OO\), the ordering of
the embeddings, and the chosen system of maximal-order ideal
representatives, which we suppress from the notation.  In fact replacing
\(I\) by \(aI\) and transporting \(w\) by \(\bigwedge^2a\) multiplies
that class's products by \(\operatorname{sgn}(\Nm_{K/\Q}a)\), since
\(\Nm(aI)=|\Nm_{K/\Q}a|\Nm I\).  We keep the representative system fixed.
A term of positive
weight admits a local label at every prime.  These labels patch to a
global \(L\), and local membership implies \(w\in\bigwedge^2_{\Z}L\).
Equation~\eqref{eq:ext-signed-products}, with positive signs, then shows
that \(\Tr z(I,w)\) is the integral value of \(Q\) on a unimodular
exterior lattice.  Thus
\begin{equation}
 M(z)>0\quad\Longrightarrow\quad
 z\in F^\times,\qquad \Tr_{F/\Q}z\in\Z.
 \label{eq:ext-trace-support}
\end{equation}

For \(m\in\Z\setminus\{0\}\), put
\[
 \mathcal H_m=\{(y_b)_{b\in\mathcal B}\in\R^{\mathcal B}:
                      \textstyle\sum_b y_b=m\}.
\]
A rectangle in \(\mathcal H_m\) means the inverse image of a product of
two intervals under any two distinct partition coordinates.  Its area
is the product of those interval lengths.  The following estimate is
the arithmetic input needed to exclude the obstruction.

\begin{theorem}[Planar exterior count]
\label{thm:exterior-count}
There is an absolute \(\tau>0\) with the following property.  Fix
\(m\in\Z\setminus\{0\}\) and a bounded set
\(\Omega\subset\mathcal H_m\).  For every primitive totally real quartic
field \(K\), every order \(\OO\subset K\), every ordering of its real
embeddings, and every chosen system of fractional \(R\)-ideal
representatives for the ordinary classes, let \(M\) be defined by
Equation~\eqref{eq:ext-multiplicity}.
For all sufficiently large \(X\), every rectangle
\(\mathcal R\subset\Omega\) whose two side lengths are at least
\(X^{-\tau}\) satisfies
\begin{equation}
 \sum_{\substack{z\in F,\ \Tr z=m\\z\in\mathcal R}}M(z)
 \ll_{m,\Omega}\kappa\sqrt D\,\operatorname{area}(\mathcal R).
 \label{eq:ext-planar-count}
\end{equation}
The implied constant and the threshold for \(X\) are uniform in the
field, the order, the embedding ordering, the ideal representative system,
the position of the rectangle, and its side lengths and aspect ratio.
\end{theorem}

Here and below we identify \(z\in F\) with its real triple.  To see the
scale of this statement, suppose \(Q(\widehat w)=m\) and the four
coordinates outside \(12|34\) have absolute value less than \(r\).
The signed products of \(\widehat w\) then have the form
\((m-y_2-y_3,y_2,y_3)\) with \(|y_2|,|y_3|<r^2\).  Their region has area
\(O(r^4)\).  Equation~\eqref{eq:ext-signed-products} allows either trace
\(m\) or \(-m\) for \(z(I,w)\), so the theorem is used at both fixed
traces.  Sections~\ref{sec:local-fourier}--\ref{sec:sieve} prove
Theorem~\ref{thm:exterior-count}; the remainder of this section supplies
its local and algebraic inputs.

\subsection{The local product fibers}

Fix \(m\ne0\) for the rest of the section.  Choose once a sufficiently
large absolute \(p_0\), and put
\begin{equation}
 S=\{p:p\mid Dqm\}\ \cup\ \{p:p\le p_0\}.
 \label{eq:ext-prime-set}
\end{equation}
For fixed \(m\),
\(\prod_{p\in S}p\ll_m Dq\le X^2\).  At primes in \(S\) we shall keep
the exact additive density of the products.  Outside \(S\), integral
valuations give a simpler bound.

The following comparison explains what a fiber remembers about an exterior
vector.  For \(k=\Q\) or \(k=\Q_p\), put \(K_k=K\otimes_\Q k\) and
\(V_k=V\otimes_\Q k\), and take embedding coordinates over a splitting
field of the etale algebra \(K_k\).

\begin{lemma}[Comparison of two product fibers]
\label{lem:ext-comparison}
Let \(w,w'\in V_k\) have all six embedding coordinates nonzero.  Suppose
that their coordinate ratios \(r_{ij}=w'_{ij}/w_{ij}\) satisfy
\[
 r_{12}r_{34}=r_{13}r_{24}=r_{14}r_{23}=s\in k^\times.
\]
There exist \(a\in K_k^\times\) and \(c\in k^\times\) such that
\[
 w'=c(\bigwedge\nolimits^2 a)w,\qquad c^2\Nm_{K_k/k}(a)=s.
\]
The pair \((a,c)\) is unique up to
\((a,c)\mapsto(da,d^{-2}c)\), \(d\in k^\times\).
\end{lemma}

\begin{proof}
Over the splitting field set
\[
 c_0=\frac{r_{12}r_{13}}{r_{23}},\qquad
 (a_1,a_2,a_3,a_4)=
 \left(1,\frac{r_{23}}{r_{13}},\frac{r_{23}}{r_{12}},
                 \frac{r_{14}r_{23}}{r_{12}r_{13}}\right).
\]
The three equal opposite products give \(r_{ij}=c_0a_i a_j\) for every
pair.  If two such factorizations are divided, all products of two of
the four coordinatewise quotients are equal.  Comparing pairs that
share one index shows that the four quotients are equal.  Hence the
projective tuple \([a_1:a_2:a_3:a_4]\) is uniquely determined.

The ratios \(r_{ij}\) are Galois-equivariant without alternating signs.
Uniqueness makes the projective tuple equivariant for the permutation
action on the four embeddings.  Express it in a \(k\)-basis of \(K_k\).
The resulting projective coefficient vector is fixed by the usual
Galois action.  Dividing by one nonzero coefficient makes every
coefficient fixed, hence in \(k\).  It gives \(a\in K_k^\times\), since
all four embedding components are nonzero.  The remaining coefficient
\(c\) in \(r_{ij}=c a_i a_j\) is Galois-fixed as well.  This proves
existence over \(k\), including when \(K_k\) is a product of nonsplit
fields.  The projective uniqueness proves the stated ambiguity.
Multiplying two opposite equalities gives \(c^2\Nm(a)=s\).
\end{proof}

For a prime \(p\), \(F_p^\times\) denotes the elements nonzero in every
field factor of \(F_p\).  In the local formulas below, \(\Nm(a)\) denotes
the algebraic norm \(\Nm_{K_p/\Q_p}(a)\).  For a global pair in
Equation~\eqref{eq:ext-multiplicity}, choose a representative
\(0\ne w\in\bigwedge^2_\Z I\) of its \(U^+\)-orbit and put
\[
 x=(\bigwedge\nolimits^2\xi_p)^{-1}w,\qquad
 u=\frac{\Nm_{K_p/\Q_p}(\xi_p)}{\Nm I}.
\]
The valuations of the numerator and denominator in \(u\) are equal:
both are the norm valuation of \(I_p\).  Thus \(u\in\Z_p^\times\), and
Equation~\eqref{eq:ext-covariance} gives \(u\Psi(x)=z(I,w)\).
Inverting the unit variable in Equation~\eqref{eq:ext-label-probability}
gives
\begin{equation}
 P_p(I,w)=
 \Pr_{a\in R_p^\times}\bigl[(\bigwedge\nolimits^2a)x\in L_{V,p}\bigr].
 \label{eq:ext-unit-membership}
\end{equation}

These representatives lead us to consider, for arbitrary \(z\in F_p^\times\),
the full local fiber
\[
 \mathscr F_p(z)=
 \{(x,u)\in V_p\times\Z_p^\times:u\Psi(x)=z\}.
\]
Define the abelian local group and its unit subgroup by
\begin{align}
 \mathscr G_p
  &=\{(a,c)\in K_p^\times\times\Q_p^\times:
                 c^2\Nm(a)\in\Z_p^\times\}/
        \{(d,d^{-2}):d\in\Q_p^\times\}, \notag\\
 \mathscr U_p
  &=\operatorname{image}(R_p^\times\times\Z_p^\times)
       \ \subset\mathscr G_p.
 \label{eq:ext-local-group}
\end{align}
The quotient in the first line means that \(d\) acts by
\((a,c)\mapsto(da,d^{-2}c)\).  The action on the fiber is
\begin{equation}
 (a,c)\cdot(x,u)=
 \left(c(\bigwedge\nolimits^2a)x,\,
                    \frac{u}{c^2\Nm(a)}\right).
 \label{eq:ext-fiber-action}
\end{equation}
Lemma~\ref{lem:ext-comparison}, with \(s=u/u'\), shows that this action
is free and transitive when the fiber is nonempty.  The explicit ratios
in its proof also give analytic local inverses, so this is an analytic
identification of the fiber with \(\mathscr G_p\).  The group
\(\mathscr U_p\) is compact and open.  The kernel of its unit
parametrization is exactly \(d\in\Z_p^\times\): if both \(a\) and \(da\)
are \(R_p\)-units, then \(v_p(d)=0\).  In particular this assertion
holds at ramified primes.

Choose a reference in a nonempty fiber and use it to identify the fiber
with \(\mathscr G_p\).  For a coset
\(\mathfrak c\in\mathscr G_p/\mathscr U_p\), let
\(\alpha_p(\mathfrak c;z)\) be the fraction of that coset for which
\(x\in L_{V,p}\), using Haar probability on \(\mathscr U_p\).
We call the coset accepted when \(\alpha_p(\mathfrak c;z)>0\).
The fraction is independent of the representative of the coset.  The sum of
these fractions is independent of the reference.  For an empty fiber
we declare this sum to be zero.

A scalar \(c\in\Z_p^\times\) preserves \(L_{V,p}\).  The pushforward of
the product Haar probability on \(R_p^\times\times\Z_p^\times\) is
Haar probability on \(\mathscr U_p\).  Hence the representative constructed
from \((I,w)\) occupies a coset \(\mathfrak c\) for which
\(P_p(I,w)=\alpha_p(\mathfrak c;z(I,w))\).

We now measure the sum over all local cosets.  Fix an additive character
\(\psi_p\) of \(\Q_p\) with conductor \(\Z_p\), and let \(dz\) be the
self-dual Haar measure on \(F_p\) for the trace pairing.  Thus
\begin{equation}
 \vol_{dz}(R_{F,p})=|D_F|_p^{1/2}.
 \label{eq:ext-self-dual-measure}
\end{equation}

For every \(p\), consider the law of the random variable
\begin{equation}
 u\Psi(x),\qquad x\ \text{uniform additive in }L_{V,p},\qquad
 u\ \text{uniform in }\Z_p^\times,
 \label{eq:ext-density-law}
\end{equation}
with the two choices independent.  The next lemma proves that this law
has a density with respect to \(dz\); we denote its locally constant
version on \(F_p^\times\) by \(f_p^*(z)\).  Only the densities for
\(p\in S\) will enter the multiplicity bound.

\begin{lemma}[Local density and disintegration]
\label{lem:ext-density}
The law in Equation~\eqref{eq:ext-density-law} has a density, with a
locally constant version on \(F_p^\times\).  For every \(z\in F_p^\times\),
only finitely many cosets are accepted, and their fractions satisfy the
exact identity
\begin{equation}
 \sum_{\mathfrak c\in\mathscr G_p/\mathscr U_p}
        \alpha_p(\mathfrak c;z)
 =
 |D|_p^{-1/2}
 \frac{1-p^{-1}}
      {\displaystyle\prod_{\mathfrak p\mid p,\ \mathfrak p\subset R}
                      (1-\Nm(\mathfrak p)^{-1})}
 f_p^*(z).
 \label{eq:ext-density-disintegration}
\end{equation}
The sum and the density are both zero if \(\mathscr F_p(z)\) is empty.
The identity holds at every prime, including primes dividing \(q\) and
the small ramified primes.
\end{lemma}

\begin{proof}
Let \(\Phi_p(x,u)=u\Psi(x)\).  If \(\Phi_p(x,u)\in F_p^\times\), all six
split coordinates of \(x\) are nonzero.  After splitting, the derivative
in \(x\) is the derivative of three products on three disjoint
two-dimensional planes.  It has rank three, and therefore has rank
three over \(\Q_p\).  The map is a submersion there.  The remaining
input set is a finite union of proper algebraic zero loci and has
additive measure zero.  Local change of variables gives an absolutely
continuous pushforward on the submersion locus, hence a density for
the full law.  To obtain a locally constant version at a fixed
\(z\in F_p^\times\), cover its inverse image in the compact set
\(L_{V,p}\times\Z_p^\times\) by finitely many submersion charts,
refined to disjoint compact open pieces.  The compact complement has
image avoiding a sufficiently small neighborhood of \(z\).
On that neighborhood each chart contributes a locally constant
density.  This also treats an empty inverse image.

Give \(V_p\) its additive measure for which \(L_{V,p}\) has mass one,
before restricting to that lattice.  We first compute the fiber form for
the product input measure after disintegration against \(dz\), and then
the mass of \(\mathscr U_p\) for that form.  In a \(\Q_p\) submersion chart,
the relative absolute four-form is obtained by evaluating the input top
form on four fiber basis vectors and three lifts of a target basis, then
dividing by the \(dz\)-form on that target basis.  We may compute this
identity after extension to a splitting field, still evaluating on \(\Q_p\) tangent
vectors and using the fixed extension of \(|p|_p=p^{-1}\).  We use the
same interpretation for the scalar quotient below.

Over a splitting field, the exterior embedding matrix has determinant
\(\delta_{\OO}^{\,3}\).  Consequently its measure form in the six
embedding coordinates is
\[
 |\delta_{\OO}|_p^{-3}\left|\bigwedge_{i<j}dx_{ij}\right|.
\]
Likewise \(dz\) is the absolute embedding-coordinate form
\(\left|\bigwedge_{b\in\mathcal B}dz_b\right|\).  To verify the latter
normalization, evaluate this form on an \(R_{F,p}\)-basis.  The square
of its embedding determinant is the discriminant, so its absolute
value is \(|D_F|_p^{1/2}\), as in
Equation~\eqref{eq:ext-self-dual-measure}.

After splitting, compute the relative form by retaining
\(x_{12},x_{13},x_{14},u\) and substituting the three product equations
for \(x_{34},x_{24},x_{23}\).  The three change-of-variable factors are
\(|\delta_{\OO}/(u x_{12})|_p\),
\(|\delta_{\OO}/(u x_{13})|_p\), and
\(|\delta_{\OO}/(u x_{14})|_p\).  They cancel
\(|\delta_{\OO}|_p^{-3}\), since \(u\) is a unit.  Haar probability on
\(\Z_p^\times\) is \((1-p^{-1})^{-1}|du/u|\).  The resulting fiber form is
\begin{equation}
 \frac{1}{1-p^{-1}}
 \left|\frac{dx_{12}}{x_{12}}\wedge
       \frac{dx_{13}}{x_{13}}\wedge
       \frac{dx_{14}}{x_{14}}\wedge\frac{du}{u}\right|.
 \label{eq:ext-fiber-form}
\end{equation}
The fiber action preserves this measure.  Indeed its determinant on
\(V_p\) is \(c^6\Nm(a)^3=(c^2\Nm(a))^3\), of absolute value one,
and it multiplies \(u\) by a unit while fixing \(z\).

It remains to compute the mass of \(\mathscr U_p\) in
Equation~\eqref{eq:ext-fiber-form}.  Put
\(\ell_i=da_i/a_i\) and \(s=dc/c\) over the splitting field.  The
pullbacks of its four logarithmic differentials under
Equation~\eqref{eq:ext-fiber-action} are
\[
 \ell_1+\ell_2+s,\quad \ell_1+\ell_3+s,\quad
 \ell_1+\ell_4+s,\quad -\ell_1-\ell_2-\ell_3-\ell_4-2s.
\]
The vertical scalar direction \(d\) has vector
\((1,1,1,1,-2)\).  Taking the first \(\ell_i\) as its coordinate and
using the slice \(\ell_1=0\), the matrix on the four remaining
differentials is
\[
 \begin{pmatrix}
  1&0&0&1\\
  0&1&0&1\\
  0&0&1&1\\
  -1&-1&-1&-2
 \end{pmatrix},
 \qquad \det=1.
\]
The vertical coordinate has coefficient \(1\).  Equivalently, the
wedge of the four displayed pullbacks is the contraction of
\((\bigwedge_i\ell_i)\wedge s\) by the vertical vector, up to sign.
Evaluating this identity on four \(\Q_p\) tangent lifts of a quotient basis
gives the same absolute quotient form for nonsplit \(K_p\).  On a patch
where a \(\Q_p\)-linear coefficient \(\lambda(a)\) is nonzero, normalizing
it to its value at a reference point gives a genuine \(\Q_p\) section;
after shrinking a unit-group patch, the normalizing scalar is a unit.
Since \(\lambda(da)=d\lambda(a)\), the vertical logarithmic derivative is
\(1\).  Thus the contraction is the quotient of the parameter form by
\(|dd/d|\), with quotient Jacobian \(1\), including at \(p=2\).

For an \(R_p\)-basis, the multiplicative embedding form on \(K_p^\times\)
is the additive basis form multiplied by its embedding determinant
and divided by \(\Nm(a)\).  On \(R_p^\times\) the latter is a unit.
The mass of this unit set is consequently
\[
 |D|_p^{1/2}
 \prod_{\mathfrak p\mid p,\ \mathfrak p\subset R}
               (1-\Nm(\mathfrak p)^{-1}).
\]
The \(c\)-unit mass and the scalar-kernel \(d\)-unit mass are both
\(1-p^{-1}\), and cancel when taking the quotient.  Haar integration
on the compact unit groups assembles the local sections.  Including the
factor \((1-p^{-1})^{-1}\) in Equation~\eqref{eq:ext-fiber-form}, the
mass of each \(\mathscr U_p\)-coset is
\begin{equation}
 \frac{|D|_p^{1/2}}{1-p^{-1}}
 \prod_{\mathfrak p\mid p,\ \mathfrak p\subset R}
               (1-\Nm(\mathfrak p)^{-1}).
 \label{eq:ext-unit-coset-mass}
\end{equation}
This calculation uses the actual image of the unit group, so it
requires no surjectivity assertion about the norm on units.

The accepted part of the fiber is the inverse image of \(z\) in the
compact set \(L_{V,p}\times\Z_p^\times\).  Since \(\mathscr U_p\) is
open, this compact set meets only finitely many of its cosets.
Integrating its indicator in the fiber multiplies the mass in
Equation~\eqref{eq:ext-unit-coset-mass} by
\(\sum_{\mathfrak c}\alpha_p(\mathfrak c;z)\).  The result is the
density \(f_p^*(z)\).  Solving for the sum proves
Equation~\eqref{eq:ext-density-disintegration}.
\end{proof}

In particular the order index has disappeared from the prefactor in
Equation~\eqref{eq:ext-density-disintegration}: its contribution through
\(\delta_{\OO}\) cancels in the three products.  Its effect remains
inside the density of the lattice \(L_{V,p}\).  The discriminant in the
prefactor is \(D\) because the translating parameters are the
maximal-order units \(R_p^\times\).

\subsection{Global multiplicity and the unramified factors}

The local identity counts every local coset.  We next bound how many
global pairs can occupy one assignment of those cosets.

\begin{lemma}[Global pairs in one assignment of cosets]
\label{lem:ext-global-cosets}
Fix \(z\in F^\times\).  Choose a reference in each nonempty local fiber
\(\mathscr F_p(z)\).  For any assignment
\(\mathfrak c_p\in\mathscr G_p/\mathscr U_p\), at most \(32\) pairs
\[
 I,\qquad w\in(\bigwedge^2_{\Z}I\setminus\{0\})/U^+,\qquad z(I,w)=z
\]
occupy the assigned coset at every prime.
\end{lemma}

\begin{proof}
For two such pairs, Lemma~\ref{lem:ext-comparison} over \(\Q\) gives
\(w'=c(\bigwedge^2a)w\) with \(a\in K^\times\), \(c\in\Q^\times\), since
the common opposite ratio is \(\Nm(I')/\Nm(I)\).  Relative to the local
generators \(\xi_p,\xi'_p\), their local comparison is represented by
\((a\xi_p/\xi'_p,c)\).  Equality of the local cosets means that some
\(d_p\in\Q_p^\times\) satisfies
\[
 d_pa\xi_p/\xi'_p\in R_p^\times,\qquad d_p^{-2}c\in\Z_p^\times.
\]
In particular \(v_p(c)=2v_p(d_p)\) is even at every finite prime.
Therefore \(c=\epsilon s^2\) with \(s\in\Q^\times\) and
\(\epsilon\in\{1,-1\}\).  Absorbing \(s\) into \(a\) makes
\(a\xi_p/\xi'_p\) a unit for every \(p\).  Hence \(I'=aI\) as fractional
maximal-order ideals.  The chosen representatives of the ordinary ideal
classes then give \(I'=I\) and \(a\in R^\times=U\).

There are at most \([U:U^+]\le16\) possibilities for the unit action
modulo \(U^+\), distinguished by the four real signs, and two
possibilities for \(\epsilon\).  This proves the bound \(32\).
\end{proof}

At \(p\notin S\), both local algebras are unramified,
\(\OO_p=R_p\), and \(\delta_{\OO}\) is a unit in a splitting field.
Let \(\varphi_p\) be Frobenius acting on the four embeddings and on
\(\mathcal B\).  Its orbits on \(\mathcal B\) correspond to primes
\(\mathfrak q\mid p\) of \(F\), of residue degrees
\(f_{\mathfrak q}\).  The map from two-element subsets to partitions
sends \(ij\) to \(ij|kl\).  Above a Frobenius orbit of partitions its
pair labels form either two orbits or one.  Define
\[
 \epsilon_{\mathfrak q}=
 \begin{cases}
  1,&\text{if the pair labels form two orbits},\\
 -1,&\text{if they form one orbit},
 \end{cases}
 \qquad
 b_{\mathfrak q}(j)=
 \begin{cases}
  j+1,&\epsilon_{\mathfrak q}=1,\\
  \mathbf1_{\{2\mid j\}},&\epsilon_{\mathfrak q}=-1
 \end{cases}
 =\sum_{r=0}^j\epsilon_{\mathfrak q}^{\,r}.
\]
For \(z\in R_{F,p}\cap F_p^\times\), put
\begin{equation}
 j_{\mathfrak q}=v_{\mathfrak q}(z),\qquad
 b_p(z)=\prod_{\mathfrak q\mid p}b_{\mathfrak q}(j_{\mathfrak q}),
 \qquad n_p(z)=\max_{\mathfrak q\mid p}j_{\mathfrak q}.
 \label{eq:ext-unramified-factor}
\end{equation}
Here \(v_{\mathfrak q}\) is the integral uniformizer valuation; it has
\(v_{\mathfrak q}(p)=1\) at these primes.  For probability statements
one may set \(n_p=\infty\), \(b_p=0\) when a component vanishes, a set
of Haar measure zero.

\begin{lemma}[Unramified local data]
\label{lem:ext-unramified-data}
Let \(p\notin S\).  A local fiber with an accepted coset has
\(z\in R_{F,p}\), and for \(z\in R_{F,p}\cap F_p^\times\),
\begin{equation}
 \sum_{\mathfrak c\in\mathscr G_p/\mathscr U_p}
       \alpha_p(\mathfrak c;z)\le b_p(z).
 \label{eq:ext-unramified-cosets}
\end{equation}
One has \(b_p=1\) when \(n_p=0\), and \(0\le b_p\le(1+n_p)^3\).

Give the affine lattice
\[
 \{z\in R_{F,p}:\Tr_{F_p/\Q_p}z=m\}
\]
its Haar probability.  With \(t=p^{-1}\) and
\(f_{\mathfrak p}\) the residue degrees of the primes of \(K\) above \(p\),
\begin{align}
 \E_p b_p
   &=(1-t)(1-t^3)
       \prod_{\mathfrak p\mid p}(1-t^{f_{\mathfrak p}})^{-1}
     =1+\frac{a_K(p)-1}{p}+O(p^{-2}), \label{eq:ext-unramified-mean}\\
 \P_p(n_p\ge k)&\le3p^{-k}\qquad(k\ge1). \label{eq:ext-unramified-tail}
\end{align}
The error constant is absolute.  In any integral affine coordinates
on the trace lattice, \(\{n_p\ge k\}\) is determined modulo \(p^k\),
and \(b_p\mathbf1_{\{n_p=k\}}\) is determined modulo \(p^{k+1}\)
for every finite \(k\ge0\).
\end{lemma}

\begin{proof}
An accepted coset has a representative \(x\in L_{V,p}=\bigwedge^2 R_p\).
Use the extended valuation with \(v_p(p)=1\) on a splitting field.
Its six valuations \(k_{ij}=v_p(x_{ij})\) are nonnegative integers,
invariant under Frobenius.  Since \(u\) and \(\delta_{\OO}\) are units,
\[
 k_{ij}+k_{kl}=j_{\mathfrak q}
 \quad\text{for the partition }ij|kl\text{ over }\mathfrak q .
\]
This proves integrality of \(z\).  If the pair labels over
\(\mathfrak q\) have two Frobenius orbits, one of their two valuations
can be chosen from \(0,\ldots,j_{\mathfrak q}\).  If there is one orbit,
the two valuations are equal, which requires even \(j_{\mathfrak q}\).
The number of records is therefore at most \(b_p(z)\).

Two representatives with the same six valuations have comparison
\((a,c)\) satisfying
\[
 v_p(c)+v_p(a_i)+v_p(a_j)=0\qquad(i<j).
\]
Subtraction makes all four \(v_p(a_i)\) equal to a common \(r\), and
\(v_p(c)=-2r\).  Because \(K_p\) is unramified, \(r\in\Z\).
Replacing \(a\) by \(p^{-r}a\) and \(c\) by \(p^{2r}c\) puts both in
the unit subgroup.  Thus a record specifies at most one coset, also
for nonsplit unramified factors.  The subgroup \(\mathscr U_p\)
preserves \(\bigwedge^2R_p\), so its coset fraction is either \(0\)
or \(1\).  This proves Equation~\eqref{eq:ext-unramified-cosets}.
The bounds at depth zero and at general depth follow immediately
from the definition and the at most three factors.

The trace from an unramified integer ring to \(\Z_p\) is onto, as is
seen after reduction from the nonzero finite-field trace.  If
\(F_p\) has more than one factor, the trace-\(m\) coset therefore
projects uniformly onto each chosen factor: the other factors can
supply any required trace.  More generally, for integers
\(r_{\mathfrak q}\ge0\), the probability that
\(v_{\mathfrak q}(z)\ge r_{\mathfrak q}\) for all \(\mathfrak q\) is
\[
 \begin{cases}
  t^{\sum_{\mathfrak q}f_{\mathfrak q}r_{\mathfrak q}},
       &\min_{\mathfrak q}r_{\mathfrak q}=0,\\
  0,&\min_{\mathfrak q}r_{\mathfrak q}>0.
 \end{cases}
\]
The first line follows by projecting onto the constrained factors
while a factor with \(r_{\mathfrak q}=0\) supplies the trace.
The second uses \(p\nmid m\).  It also covers the case that \(F_p\)
is a field.

Expanding \(b_{\mathfrak q}(j)=\sum_{r=0}^j\epsilon_{\mathfrak q}^r\)
and summing the absolutely convergent geometric series gives
\begin{equation}
 \E_p b_p=
 \frac{1-(\prod_{\mathfrak q}\epsilon_{\mathfrak q})t^3}
      {\prod_{\mathfrak q}(1-\epsilon_{\mathfrak q}t^{f_{\mathfrak q}})}.
 \label{eq:ext-geometric-mean}
\end{equation}
The five Frobenius types give the following data.  A pair
\((f,\epsilon)\) in the second column is attached to one prime of \(F\).
\[
\begin{array}{c|c|c|c}
 \text{cycle type on four labels}
 & (f_{\mathfrak q},\epsilon_{\mathfrak q})
 & \prod_{\mathfrak p\mid p}(1-t^{f_{\mathfrak p}})
 & a_K(p)-1\\ \hline
 1^4 &(1,+),(1,+),(1,+)&(1-t)^4&3\\
 2\,1^2 &(1,+),(2,+)&(1-t)^2(1-t^2)&1\\
 2^2 &(1,+),(1,-),(1,-)&(1-t^2)^2&-1\\
 3\,1 &(3,+)&(1-t)(1-t^3)&0\\
 4 &(1,-),(2,-)&1-t^4&-1
\end{array}
\]
In every row \(\prod_{\mathfrak q}\epsilon_{\mathfrak q}=1\) and
\[
 \prod_{\mathfrak q}(1-\epsilon_{\mathfrak q}t^{f_{\mathfrak q}})
   =\frac{\prod_{\mathfrak p\mid p}(1-t^{f_{\mathfrak p}})}{1-t}.
\]
Substitution in Equation~\eqref{eq:ext-geometric-mean} proves the exact
mean.  Its linear term is the number of degree-one primes of \(K\)
minus one, which is \(a_K(p)-1\); the five fixed rows give the uniform
error term.

If \(F_p\) is a field, \(n_p\ge1\) would make its trace divisible by
\(p\), and is impossible.  Otherwise the uniform projection above
gives
\(\P_p(v_{\mathfrak q}(z)\ge k)=p^{-f_{\mathfrak q}k}\).
Taking the union over at most three factors proves
Equation~\eqref{eq:ext-unramified-tail}.  Finally, the trace map is
onto, so the affine trace lattice has two integral coordinates.
Projection to each unramified factor modulo \(p^{k+1}\) is determined
by these coordinates modulo \(p^{k+1}\).  The valuations truncated
at \(k+1\), and hence the asserted threshold and finite-depth
functions, have the stated congruence dependence.  This remains true
on the component-vanishing set with the convention above.
\end{proof}

\begin{proposition}[Multiplicities from local densities]
\label{prop:ext-multiplicity-majorant}
For \(z\in F\) with \(\Tr z=m\ne0\), a positive value of \(M(z)\)
requires \(z\) to be integral outside \(S\).  For such integral \(z\),
\begin{equation}
 M(z)\le32
 \prod_{p\in S}
 \left(
 |D|_p^{-1/2}
 \frac{1-p^{-1}}
      {\displaystyle\prod_{\mathfrak p\mid p,\ \mathfrak p\subset R}
                         (1-\Nm(\mathfrak p)^{-1})}
 f_p^*(z)\right)
 \prod_{p\notin S}b_p(z).
 \label{eq:ext-local-majorant}
\end{equation}
All products in this assertion are finite at every contributing \(z\).
\end{proposition}

\begin{proof}
The field \(F\) and the condition \(\Tr z=m\ne0\) imply
\(z\in F^\times\).  At every prime the set of accepted cosets is
finite by Lemma~\ref{lem:ext-density}.  At \(p\notin S\),
Lemma~\ref{lem:ext-unramified-data} proves the integrality requirement
and bounds the local sum by \(b_p(z)\).  The set
\[
 E_z=S\cup\{p:z\notin R_{F,p}^\times\}
\]
is finite, since a nonzero global \(z\) is a unit at all but finitely
many primes.  At \(p\notin E_z\), \(b_p=1\), and each coset fraction is
\(0\) or \(1\), so there is at most one accepted coset, of fraction
\(1\).  If some prime has no accepted coset, then \(M(z)=0\) and the
claimed inequality is immediate.  Otherwise the accepted coset is
unique with fraction \(1\) outside \(E_z\), and each prime inside the
finite set \(E_z\) has finitely many accepted cosets.  Thus there are
only finitely many assignments of accepted cosets, and all local sums
outside \(E_z\) equal \(1\).

Choose local references.  Equation~\eqref{eq:ext-unit-membership}
identifies the weight of a global pair at each prime with its coset
fraction \(\alpha_p\).  By Lemma~\ref{lem:ext-global-cosets}, at most
\(32\) global pairs occupy any assignment of accepted cosets.
Summing over the finite set of assignments, including ones not
realized globally, bounds \(M(z)\) by \(32\) times the product of the
local sums of \(\alpha_p\).  Use
Equation~\eqref{eq:ext-density-disintegration} at \(p\in S\) and
Equation~\eqref{eq:ext-unramified-cosets} outside \(S\).
This proves Equation~\eqref{eq:ext-local-majorant} and finiteness.
\end{proof}

The unramified mean in Equation~\eqref{eq:ext-unramified-mean} has the
Euler factor needed for the residue \(\kappa\).  The remaining problem
in Theorem~\ref{thm:exterior-count} is to average the densities
\(f_p^*\) for \(p\in S\) before inserting these unramified factors.
For a frequency \(\alpha\in F_p\), testing the law \(u\Psi(x)\) against
the character \(z\mapsto\psi_p(\Tr_{F_p/\Q_p}(\alpha z))\) produces the
quadratic phase \(u\Tr_{F_p/\Q_p}(\alpha\Psi(x))\).
We express these forms through a single unimodular bilinear form and an
action of \(F\).  The elements of \(F\) whose operators preserve \(L_V\)
will then form the order that controls the local Fourier analysis.

\subsection{The integral quadratic structure and its order}

Let \(B_V\) be the polarization of \(Q_V\) without division by two:
\[
 B_V(x,y)=Q_V(x+y)-Q_V(x)-Q_V(y).
\]
Over a splitting field, write \(V_b\) for the plane spanned by the
two coordinate axes belonging to the partition \(b\), the same
coordinate planes used over \(\R\) in Section~\ref{sec:rigidity}.
For an equivariant triple \(\alpha=(\alpha_b)\in F\), let \(\alpha_b\)
scale both coordinate axes of \(V_b\).

\begin{lemma}[Self-adjoint quadratic structure]
\label{lem:ext-quadratic-structure}
The plane-scalar action descends to a rational action
\(\rho:F\to\End_{\Q}(V)\) for which \(V\) has rank two over \(F\).
The form \(Q_V\) in an exterior \(\Z\)-basis of \(L_V\) is the three
split pairs in Equation~\eqref{eq:ext-integral-form}; in particular
\(B_V\) is unimodular on \(L_V\).  The action \(\rho\) is self-adjoint
for \(B_V\), and for \(\alpha\in F\),
\begin{align}
 Q_\alpha(x):=\Tr_{F/\Q}(\alpha\Psi(x))
    &=\tfrac12 B_V(\rho(\alpha)x,x), \label{eq:ext-polar-operator}\\
 B_\alpha(x,y)&=B_V(\rho(\alpha)x,y), \notag\\
 \det_{\Q}\rho(\alpha)&=\Nm_{F/\Q}(\alpha)^2,\qquad
 \Psi(\rho(\alpha)x)=\alpha^2\Psi(x). \label{eq:ext-rank-two}
\end{align}
Here \(B_\alpha\) is the polarization of \(Q_\alpha\).
The same assertions hold after base change to every \(\Q_p\).
\end{lemma}

\begin{proof}
Galois permutes the planes \(V_b\), including possible alternating
signs on their coordinate axes, and permutes the scalars \(\alpha_b\)
in the same way.  The plane-scalar action therefore commutes with
Galois and descends.  After splitting there are three planes of
dimension two, proving the rank assertion.  Equation~\eqref{eq:ext-integral-form}
shows that the matrix of \(B_V\) pairs the three opposite pairs by
unit off-diagonal entries, with the displayed signs, so it is
unimodular.  On each plane, multiplying both coordinates by
\(\alpha_b\) is self-adjoint for its split pairing.  The contribution
of that plane to \(\frac12B_V(\rho(\alpha)x,x)\) is
\(\alpha_b\Psi_b(x)\), which proves
Equation~\eqref{eq:ext-polar-operator}.  Polarization gives the next
identity.  The determinant and the last quadratic identity follow
by multiplying, respectively, the two equal eigenvalues and the
two coordinate scalars on each plane.  All these are rational
identities, hence persist under base change.
\end{proof}

Define the stabilizer order and its discriminant by
\begin{equation}
 C_F=\{\alpha\in F:\rho(\alpha)L_V\subset L_V\},\qquad
 \Delta_F=\Disc(C_F).
 \label{eq:ext-stabilizer-order}
\end{equation}
It is a ring containing \(1\) and a full \(\Z\)-lattice: this follows
from the rationality of \(\rho\) by clearing denominators on a basis
of \(F\).  It is contained in \(R_F\), since an \(\alpha\) whose
matrix on \(L_V\) is integral is a root of that integral monic
characteristic polynomial.  Thus it is an order.
The stabilizer condition is a finite list of matrix-entry integrality
conditions in a lattice basis, so it commutes with completion:
\begin{equation}
 C_{F,p}=C_F\otimes_\Z\Z_p
   =\{\alpha\in F_p:\rho(\alpha)L_{V,p}\subset L_{V,p}\}.
 \label{eq:ext-local-stabilizer}
\end{equation}
The following comparison ensures that its discriminant tends to infinity
with the discriminant of \(\OO\), even when the field \(K\) stays fixed.

\begin{lemma}[Scale of the stabilizer order]
\label{lem:ext-order-scale}
There is an absolute constant \(C\) such that
\begin{equation}
 \frac12 X^{1/2}\le
 \Delta_F=D_F[R_F:C_F]^2\le C X^{28}.
 \label{eq:ext-order-scale}
\end{equation}
More precisely, with \(d=\Disc(\OO)=X^2\),
\[
 d^2R_F\subset C_F\subset R_F,\qquad D_F\ll D^2.
\]
In particular \(C_{F,p}=R_{F,p}\) for \(p\nmid Dq\).
\end{lemma}

\begin{proof}
Let \(E=\bigwedge^2 A_{\OO}\) be the exterior embedding matrix.
Its entries are integral in every local splitting ring, and
\(\det E=\delta_{\OO}^{\,3}\).  Since \(d=\delta_{\OO}^{\,2}\),
\[
 d^2 E^{-1}=\delta_{\OO}\operatorname{adj}(E)
\]
also has integral entries.  For \(\alpha\in R_{F,p}\), its
plane-scalar matrix over the splitting ring is integral.
Consequently
\(d^2\rho(\alpha)=d^2E^{-1}\diag(\alpha_b,\alpha_b)E\)
is integral over the splitting ring.  Its entries
are in \(\Q_p\), so they lie in \(\Z_p\).  This proves
\(d^2R_F\subset C_F\) and \([R_F:C_F]\le d^6\).

If \(p\nmid D\), the etale algebra \(R_p\) splits over an unramified
extension of \(\Z_p\); its splitting field and cubic resolvent are
unramified.  Hence every prime dividing \(D_F\) divides \(D\).
For \(p>3\), a cubic algebra is tamely ramified and its discriminant
exponent is at most \(2\).  At \(2\) and \(3\), the finitely many
local extensions of degree at most three have bounded discriminant
exponents \cite[Proposition~7.64 and Remark~7.65]{MilneANT}.
It follows that \(D_F\ll D^2\).  The discriminant-index formula now
gives
\[
 \Delta_F=D_F[R_F:C_F]^2\ll D^2d^{12}\le X^{28}.
\]
The containment \(d^2R_F\subset C_F\) also proves equality of the
local orders away from \(Dq\).

For the reverse inequality, we first explain why a projection
denominator controls \(\Disc(\OO)\).  Write \(\operatorname{tr}\) for
matrix trace, put \(\mathscr E=\End_\Z(\OO)\), and let
\(\mathscr W_K\subset\End_\Q(K)\) be the rational subspace of
multiplication matrices.  Its intersection
\(\mathscr L=\mathscr W_K\cap\mathscr E\) is saturated.  It is
multiplication by \(\OO\): if \(\alpha\OO\subset\OO\), then
\(\alpha=\alpha\cdot1\in\OO\), and the converse follows because
\(\OO\) is a ring.  The pairing
\((T,S)\mapsto\operatorname{tr}(TS)\) is unimodular on
\(\mathscr E\), since its matrix units are dual.  Its restriction
to \(\mathscr W_K\) is the nondegenerate field trace pairing.
Let \(\mathcal E_K\) be the orthogonal projection onto
\(\mathscr W_K\).

For the restricted pairing one has
\begin{equation}
 \mathcal E_K(\mathscr E)=\mathscr L^\vee.
 \label{eq:ext-projected-dual}
\end{equation}
One inclusion follows by pairing a projected integral matrix with
\(\mathscr L\).  Conversely, an integer functional on the primitive
sublattice \(\mathscr L\) extends to one on \(\mathscr E\).
Unimodularity represents the extension by an element of
\(\mathscr E\); its orthogonal projection represents the original
functional on \(\mathscr L\).  Nondegeneracy identifies that
projection with the prescribed element of \(\mathscr L^\vee\).
Consequently, if a positive integer \(N\) satisfies
\(N\mathcal E_K(\mathscr E)\subset\mathscr E\), then
\[
 \mathscr L^\vee\subset N^{-1}\mathscr L,\qquad
 X^2=\Disc(\OO)=[\mathscr L^\vee:\mathscr L]\le N^4.
\]
It remains to produce the denominator \(N=2\Delta_F\).

Choose trace-dual bases \(a_i\) of \(C_F\) and \(b_i\) of \(C_F^\vee\),
where
\[
 C_F^\vee=\{\beta\in F:\Tr_{F/\Q}(\beta C_F)\subset\Z\}.
\]
The adjugate of the trace Gram matrix gives
\(\Delta_F b_i\in C_F\).  Define on \(\End_\Q(V)\)
\begin{equation}
 \mathcal E_F(T)=\sum_{i=1}^3\rho(a_i)T\rho(b_i).
 \label{eq:ext-centralizer-projection}
\end{equation}
For two embeddings \(\sigma,\tau\) of \(F\), the dual-basis identity is
\[
 \sum_i\sigma(a_i)\tau(b_i)=
 \begin{cases}1,&\sigma=\tau,\\0,&\sigma\ne\tau.\end{cases}
\]
Thus, after splitting, \(\mathcal E_F\) extracts exactly the three
diagonal blocks \(\End(V_b)\).  It is the orthogonal projection
for the matrix trace pairing, and
\begin{equation}
 \Delta_F\mathcal E_F(\End_\Z(L_V))\subset\End_\Z(L_V).
 \label{eq:ext-centralizer-denominator}
\end{equation}
This denominator conclusion uses the trace-dual lattice only; it
does not require that lattice to be an invertible \(C_F\)-ideal.

We transfer this projection to matrices on \(K\) explicitly.
Let
\[
 j(T)(v\wedge w)=Tv\wedge w+v\wedge Tw
 \qquad(T\in\End_\Q(K)).
\]
There is a natural contraction
\(\mathfrak c:\End_\Q(\bigwedge^2K)\to\End_\Q(K)\).
For clarity, in a basis \(e_1,\ldots,e_4\) write
\[
 S(e_j\wedge e_l)=\sum_{i<k}S^{ik}_{jl}e_i\wedge e_k\qquad(j<l),
\]
and extend these coefficients antisymmetrically in \(i,k\) and in
\(j,l\), with zero coefficients when either pair repeats.  Then
\[
 \mathfrak c(S)^i_j=\sum_{k=1}^4 S^{ik}_{jk}.
\]
This is the contraction of
\(\bigwedge^2K\otimes\bigwedge^2K^\vee\) in one covariant and one
contravariant index, so it is independent of the basis.  Both \(j\)
and \(\mathfrak c\) are integral in the \(\OO\) and exterior
\(\OO\)-bases.  Expanding their coefficients gives
\begin{equation}
 \mathfrak c(j(T))=2T+\operatorname{tr}(T)\id_K.
 \label{eq:ext-contraction}
\end{equation}
Indeed an off-diagonal coefficient of \(T\) occurs for the two
remaining indices \(k\); a diagonal coefficient gives
\(\sum_{k\ne i}(T^i_i+T^k_k)=2T^i_i+\operatorname{tr}(T)\).

In embedding coordinates, an off-diagonal entry of \(j(T)\) changes
one index of a two-element subset.  Such subsets belong to distinct
opposite-pair planes.  It follows that
\(\mathcal E_F(j(T))=j(\diag T)\) over the splitting field.
Equation~\eqref{eq:ext-contraction} now proves the formula
\begin{equation}
 \mathcal E_K(T)=
 \tfrac12\bigl(\mathfrak c(\mathcal E_F(j(T)))-\operatorname{tr}(T)\id_K\bigr)
 \label{eq:ext-multiplication-projection}
\end{equation}
for the projection defined above: its right side is \(\diag T\)
over the splitting field.  A Galois-equivariant diagonal tuple is
the tuple of embeddings of an element of \(K\), so these diagonal
matrices are precisely \(\mathscr W_K\).  Diagonal extraction is
orthogonal for matrix trace.  Finally,
Equation~\eqref{eq:ext-centralizer-denominator} and integrality of
\(j\) and \(\mathfrak c\) give
\[
 2\Delta_F\mathcal E_K(\End_\Z(\OO))\subset\End_\Z(\OO).
\]
Taking \(N=2\Delta_F\) in the preceding lattice bound gives
\[
 X^2\le(2\Delta_F)^4.
\]
This is the lower inequality in Equation~\eqref{eq:ext-order-scale}.
\end{proof}

For \(x\in L_{V,p}\) and \(\alpha\in C_{F,p}\),
Equation~\eqref{eq:ext-polar-operator} gives two integral pairings:
\[
 \Tr\Psi(x)=Q_V(x)\in\Z_p,\qquad
 2\Tr(\alpha\Psi(x))=B_V(\rho(\alpha)x,x)\in\Z_p.
\]
They remain integral after multiplying \(\Psi(x)\) by a \(\Z_p\)-unit.
Thus the products pair integrally with the scalar \(1\) and with
\(2C_{F,p}\).  Define the order and its trace dual
\[
 C_F^\sharp=\Z+2C_F,\qquad
 (C_F^\sharp)^\vee
   =\{z\in F:\Tr_{F/\Q}(zC_F^\sharp)\subset\Z\}.
\]
This trace dual gives the integral support for the additive count.
The same support will bound the total depth of the unramified factors
and prevent a nonzero real product coordinate from being too small.

\begin{lemma}[Trace-dual support and product size]
\label{lem:ext-support}
One has
\begin{equation}
 [C_F:C_F^\sharp]=4,\qquad
 \Delta_F^\sharp:=\Disc(C_F^\sharp)=16\Delta_F\ll X^{28},
 \label{eq:ext-sharp-discriminant}
\end{equation}
and \(\Tr:(C_F^\sharp)^\vee\to\Z\) is onto.
For \(p\notin S\), \(C_{F,p}^\sharp=(C_{F,p}^\sharp)^\vee=R_{F,p}\).
For every prime \(p\), the density \(f_p^*\) on \(F_p^\times\) is
supported on \((C_{F,p}^\sharp)^\vee\).  Consequently
\[
 M(z)>0\ \Longrightarrow\ z\in(C_F^\sharp)^\vee.
\]
The following stronger containment is used to average the majorant in
Proposition~\ref{prop:ext-multiplicity-majorant}: if \(z\in F^\times\) is
integral outside \(S\) and \(\prod_{p\in S}f_p^*(z)>0\), then
\(z\in(C_F^\sharp)^\vee\).

Fix \(m\ne0\) and a bounded
\(\Omega\subset\mathcal H_m\).  Every
\(z\in(C_F^\sharp)^\vee\cap\Omega\) satisfies
\begin{align}
 &\Delta_F^\sharp z\in R_F,\qquad
 1\le|\Nm_{F/\Q}(\Delta_F^\sharp z)|\ll_\Omega X^{84},
       \label{eq:ext-norm-bound}\\
 &\prod_{p\notin S}p^{n_p(z)}
       \le|\Nm_{F/\Q}(\Delta_F^\sharp z)|,\qquad
 |z_b|\gg_\Omega X^{-84}\quad(b\in\mathcal B).
 \label{eq:ext-depth-coordinate-bound}
\end{align}
The constants are uniform in all arithmetic data and in the point
of \(\Omega\).
\end{lemma}

\begin{proof}
Since \(C_F\) is an order, \(C_F\cap\Q=\Z\), so \(1\) is primitive
in its lattice.  A basis \(1,a_2,a_3\) of \(C_F\) gives the basis
\(1,2a_2,2a_3\) of \(C_F^\sharp\).  This proves the index and
discriminant assertions in Equation~\eqref{eq:ext-sharp-discriminant}.
The functional taking \(1\) to \(1\) and the other two basis vectors
to \(0\) is represented in the trace dual.  Hence trace from that
dual is onto \(\Z\).  Lemma~\ref{lem:ext-order-scale}, and \(2\in S\),
give \(C_{F,p}^\sharp=R_{F,p}\) outside \(S\).

The two integrality identities preceding the lemma show that the
law in Equation~\eqref{eq:ext-density-law} is supported on
\((C_{F,p}^\sharp)^\vee\).  For a contributing global pair, use its local
representation \(z=u\Psi(x)\) and Equation~\eqref{eq:ext-fiber-action}
to translate \(x\) into \(L_{V,p}\), adjusting \(u\) by the inverse
norm of that unit.  This proves local trace-dual membership
at every prime.  Intersecting the local lattices proves the stated
global support.  For the second support assertion, the density factors
give this membership at \(p\in S\).  At \(p\notin S\), the order is
\(R_{F,p}\), whose trace pairing is unimodular because \(F_p\) is
unramified; thus \(R_{F,p}^\vee=R_{F,p}\), and the assumed integrality
gives membership there as well.

The adjugate of the trace Gram matrix of \(C_F^\sharp\) gives
\(\Delta_F^\sharp(C_F^\sharp)^\vee\subset C_F^\sharp\subset R_F\).
If \(z\in\Omega\), then \(\Tr z=m\ne0\), so \(z\ne0\); since \(F\)
is a field, all its embeddings are nonzero.  The norm
\(\Nm(\Delta_F^\sharp z)\) is consequently a nonzero integer.  The
boundedness of \(\Omega\) and \(\Delta_F^\sharp\ll X^{28}\) give
\[
 1\le|\Nm(\Delta_F^\sharp z)|
       =(\Delta_F^\sharp)^3\prod_{b\in\mathcal B}|z_b|
       \ll_\Omega X^{84}.
\]
At \(p\notin S\), \(z\) is integral and the local algebra is
unramified.  Hence
\[
 n_p(z)\le\sum_{\mathfrak q\mid p}
               f_{\mathfrak q}v_{\mathfrak q}(z)
       \le v_p\bigl(\Nm(\Delta_F^\sharp z)\bigr).
\]
The product of these inequalities gives the depth bound because
\(\Nm(\Delta_F^\sharp z)\) is an integer.  Finally, if the coordinates of
\(\Omega\) have absolute value at most \(B\), with \(B\ge1\), then
for each \(b\)
\[
 |z_b|=\frac{|\Nm(\Delta_F^\sharp z)|}
              {(\Delta_F^\sharp)^3\prod_{b'\ne b}|z_{b'}|}
       \ge\frac{1}{(\Delta_F^\sharp)^3 B^2}\gg_\Omega X^{-84}.
\]
This proves Equation~\eqref{eq:ext-depth-coordinate-bound}.
\end{proof}
